% Options for packages loaded elsewhere
\PassOptionsToPackage{unicode}{hyperref}
\PassOptionsToPackage{hyphens}{url}
\PassOptionsToPackage{dvipsnames,svgnames,x11names}{xcolor}
\documentclass[
  ngerman,
  10pt,
  a4paper,
]{article}
\usepackage{xcolor}
\usepackage[margin=22mm]{geometry}
\usepackage{amsmath,amssymb}
\setcounter{secnumdepth}{-\maxdimen} % remove section numbering
\usepackage{iftex}
\ifPDFTeX
  \usepackage[T1]{fontenc}
  \usepackage[utf8]{inputenc}
  \usepackage{textcomp} % provide euro and other symbols
\else % if luatex or xetex
  \usepackage{unicode-math} % this also loads fontspec
  \defaultfontfeatures{Scale=MatchLowercase}
  \defaultfontfeatures[\rmfamily]{Ligatures=TeX,Scale=1}
\fi
\usepackage{lmodern}
\ifPDFTeX\else
  % xetex/luatex font selection
  \setmainfont[Path=/Users/stefanhamann/.cache/codex-runtimes/codex-primary-runtime/dependencies/native/libreoffice-headless/libreoffice/LibreOfficeDev.app/Contents/Resources/fonts/truetype/,Extension=.ttf,BoldFont=DejaVuSans-Bold,ItalicFont=DejaVuSans-Oblique,BoldItalicFont=DejaVuSans-BoldOblique]{DejaVuSans}
  \setmonofont[Path=/Users/stefanhamann/.cache/codex-runtimes/codex-primary-runtime/dependencies/native/libreoffice-headless/libreoffice/LibreOfficeDev.app/Contents/Resources/fonts/truetype/,Extension=.ttf,BoldFont=DejaVuSansMono-Bold,ItalicFont=DejaVuSansMono-Oblique,BoldItalicFont=DejaVuSansMono-BoldOblique]{DejaVuSansMono}
\fi
% Use upquote if available, for straight quotes in verbatim environments
\IfFileExists{upquote.sty}{\usepackage{upquote}}{}
\IfFileExists{microtype.sty}{% use microtype if available
  \usepackage[]{microtype}
  \UseMicrotypeSet[protrusion]{basicmath} % disable protrusion for tt fonts
}{}
\makeatletter
\@ifundefined{KOMAClassName}{% if non-KOMA class
  \IfFileExists{parskip.sty}{%
    \usepackage{parskip}
  }{% else
    \setlength{\parindent}{0pt}
    \setlength{\parskip}{6pt plus 2pt minus 1pt}}
}{% if KOMA class
  \KOMAoptions{parskip=half}}
\makeatother
\ifLuaTeX
\usepackage[bidi=basic,shorthands=off]{babel}
\else
\usepackage[bidi=default,shorthands=off]{babel}
\fi
\ifPDFTeX
\else
\babelfont{rm}[Path=/Users/stefanhamann/.cache/codex-runtimes/codex-primary-runtime/dependencies/native/libreoffice-headless/libreoffice/LibreOfficeDev.app/Contents/Resources/fonts/truetype/,Extension=.ttf,BoldFont=DejaVuSans-Bold,ItalicFont=DejaVuSans-Oblique,BoldItalicFont=DejaVuSans-BoldOblique]{DejaVuSans}
\fi
\ifLuaTeX
  \usepackage{selnolig} % disable illegal ligatures
\fi
\setlength{\emergencystretch}{3em} % prevent overfull lines
\providecommand{\tightlist}{%
  \setlength{\itemsep}{0pt}\setlength{\parskip}{0pt}}
\usepackage{fvextra}
\RecustomVerbatimEnvironment{verbatim}{Verbatim}{breaklines=true,fontsize=\footnotesize}
\usepackage{xurl}
\usepackage{seqsplit}
\usepackage{adjustbox}
\usepackage{ragged2e}
\usepackage{titlesec}
\titleformat{\section}{\large\bfseries\raggedright}{}{0em}{}
\titleformat{\subsection}{\normalsize\bfseries\raggedright}{}{0em}{}
\DeclareRobustCommand{\codeword}[1]{\texttt{\seqsplit{#1}}}
\AtBeginEnvironment{longtable}{\small\RaggedRight}
\usepackage{newunicodechar}
\newunicodechar{𝒞}{\ensuremath{\mathcal C}}
\newunicodechar{𝔊}{\ensuremath{\mathfrak G}}
\newunicodechar{𝔤}{\ensuremath{\mathfrak g}}
\newunicodechar{𝓗}{\ensuremath{\mathcal H}}
\newunicodechar{ℋ}{\ensuremath{\mathcal H}}
\newunicodechar{𝔄}{\ensuremath{\mathfrak A}}
\newunicodechar{𝟙}{\ensuremath{\mathbf 1}}
\newunicodechar{𝕀}{\ensuremath{I}}
\newunicodechar{𝕋}{\ensuremath{\mathbb T}}
\newunicodechar{𝟘}{\ensuremath{\mathbf 0}}
\newunicodechar{⊞}{\ensuremath{\boxplus}}
\newunicodechar{∘}{\ensuremath{\circ}}
\newunicodechar{⟦}{\ensuremath{[\![}}
\newunicodechar{⟧}{\ensuremath{]\!]}}
\setlength{\emergencystretch}{4em}
\setlength{\parindent}{0pt}
\setlength{\parskip}{0.45em}
\AtBeginDocument{\hypersetup{colorlinks=true,linkcolor=black,urlcolor=blue}\pdfstringdefDisableCommands{\def\codeword#1{#1}}\RaggedRight}
\usepackage{fancyhdr}
\pagestyle{fancy}
\fancyhf{}
\fancyhead[L]{\small TFPT / Universalraum - v1.6.9}
\fancyfoot[R]{\thepage}
\setlength{\headheight}{15pt}
% The preserved full history now has three-digit page numbers.
\makeatletter
\renewcommand{\@pnumwidth}{2.6em}
\renewcommand{\@tocrmarg}{3.6em}
\makeatother

\usepackage{bookmark}
\IfFileExists{xurl.sty}{\usepackage{xurl}}{} % add URL line breaks if available
\urlstyle{same}
\hypersetup{
  pdflang={de-DE},
  colorlinks=true,
  linkcolor={Maroon},
  filecolor={Maroon},
  citecolor={Blue},
  urlcolor={Blue},
  pdfcreator={LaTeX via pandoc}}

\author{}
\date{}

\begin{document}

\section{TFPT / Universalraum: Update
v1.6.9}\label{tfpt-universalraum-update-v1.6.9}

\begin{enumerate}
\def\labelenumi{\arabic{enumi}.}
\setcounter{enumi}{14}
\tightlist
\item
  September 2026 · ausschließlich Änderungen gegenüber v1.6.8
\end{enumerate}

\subsection{Was jetzt neu ist}\label{was-jetzt-neu-ist}

\textbf{Die zusätzliche Zusammensetzung kann in zwei kleinen
Modellfamilien aus ihrer höheren Antwort zurückberechnet werden.}
Zugleich wurde der native kleine Eingriffsversuch vereinfacht und der
relativistische Feldanschluss bis zu einem konkreten Kinetiktest
vorangebracht. Eine vollständige TOE folgt daraus nicht.

\subsubsection{1. Ursprüngliche Quelle und globale Wiederholung sind
verschieden}\label{urspruxfcngliche-quelle-und-globale-wiederholung-sind-verschieden}

W verwendet bereits 64 gemeinsame Fermionmoden; jeder erscheint in 15
der 480 Vertices. Eine Wiederholung des ursprünglichen passiven Clocks
erzeugt keine sechs unabhängigen hellen Räume: Der Vereinigungsrang
bleibt exakt \textbf{60 statt 360}. Die alte K4-/C16-Architekturauswahl
gilt nur für ihren zusätzlichen Vertrag.

Die offene Regel betrifft die Überlappung mehrerer vollständiger
Vorkommen. Zwei überlappende Charts \(f^L=a\), \(f^R=ca+\sqrt{1-c^2}d\),
mit je ursprünglichem W und eigener Bosonbank, liefern

\[\adjustbox{max width=\linewidth}{$\displaystyle 
P(L\to R;t)=16g^4c^4t^4+O(t^6).
$}\]

Mit linker G-invarianter Bosonmischung und rechter Gesamt-Bosonzahl ist
der Vergleich vollständig innerlich invariant. Die Fälle c=3/5 und 4/5
behalten dieselben einzelnen lokalen W-Daten, haben aber verschiedene
Antworten. Der vierte Energiemoment identifiziert die Überlappung:

\[\adjustbox{max width=\linewidth}{$\displaystyle 
I_4=\frac{\mu_4-3\mu_1^2\mu_2+2\mu_1^4}{(\mu_2-\mu_1^2)^2}=1+c^4.
$}\]

Die ursprünglichen Quellstrahlen enthalten geeignete diskrete
Überlappungszahlen. Ihre Zuordnung zu diesen globalen CAR-Charts ist
jedoch noch nicht typisiert hergeleitet.

\subsubsection{2. Gesamtes Zweiladungsspektrum gleich - höhere Dynamik
verschieden}\label{gesamtes-zweiladungsspektrum-gleich---huxf6here-dynamik-verschieden}

Eine andere Familie hat eine gemeinsame Bosonbank und normierte
symmetrische Kopplung \(M_A=C\otimes A_A\). Die verteilte Wahl
\(C=I_L/\sqrt L\) und die kollektive Wahl \(C=J_L/L\) besitzen dieselben
\textbf{gesamten N=2-Spektren}. Im N=3-Sektor gibt es bei L=2 und
Energie Δ dagegen \textbf{0 beziehungsweise 64 Eigenzustände}, für
Δ\textgreater0 und g≠0. Das ist kein Port-Umbenennungseffekt.

Die allgemeine, inklusive komplexer Konvention direkt kontrollierte
Identität ist

\[\adjustbox{max width=\linewidth}{$\displaystyle 
G_C^{(3)}=8I-Q\otimes K,\qquad Q=C^\dagger C,\qquad\operatorname{tr}K=960.
$}\]

Daher folgt die positive Umkehrung

\[\adjustbox{max width=\linewidth}{$\displaystyle 
\boxed{Q=\frac1{960}\operatorname{Tr}_{\rm intern}(8I-G_C^{(3)}).}
$}\]

Benötigt wird der matrixaufgelöste Antwortoperator samt
Tensorfaktordictionary, nicht nur eine Spektralliste. Die
Anfangskrümmung der Bosonzahlantwort liefert dessen Erwartungswerte. Die
Formel rekonstruiert Q innerhalb dieser Familie; sie wählt C nicht
ursprünglich aus. Anders als im Zwei-Chart-Modell bleiben hier die
lokalen Einzelkopienamplituden nicht unverändert.

\subsubsection{3. Einfacherer Eingriff und richtiger
Recorder}\label{einfacherer-eingriff-und-richtiger-recorder}

Der neue Bosonzahlimpuls \(Z_b=e^{i\pi N_b}\) liegt auf dem gemeinsamen
Dreizustandsraum in der dokumentierten Algebra aus X und Nb; der frühere
Einzelmodenimpuls Z4 liegt außerhalb. Der extra Modenselektor entfällt.

Bei g/Δ=1/20 steigt die unbedingte n5-Wahrscheinlichkeit von
\textbf{0.0002194998817} auf \textbf{0.0005299169088}. Ein
bereitgestellter neutraler Zeiger, der kohärent nur Paar/Boson
unterscheidet, erhält alle Paarkohärenzen und liefert beim Ignorieren
des Zeigers genau den halben Folgeeffekt.

Der Impuls deponiert Δ/54, der Recorder Δ/108 im System. Separate
Schaltbarkeit, Zeigerkopplung, Anfangspräparation und terminale Auslese
bleiben Zusatzressourcen. Der Impulsarm endet anders als in v1.6.8 nicht
bosonfrei. Reine innere Symmetrieimpulse, die mit H und dem anfänglichen
Empfänger kommutieren, können den verzögerten Effekt nicht erzeugen.

\subsubsection{4. Minimaler Lorentzadapter - aber der freie Ansatz ist
null}\label{minimaler-lorentzadapter---aber-der-freie-ansatz-ist-null}

Für den gemischthändigen sechsdimensionalen Träger existiert mit einem
Impuls P eine eindeutige lineare Auslese C\_P in einen Weylspinor, Rang
zwei für P≠0. Ein Einsetzer erfüllt \(C_PJ_P=3\det(P)I\). Zeitartig gibt
es einen kovarianten Rückweg; lichtartig braucht dieser eine zusätzliche
Referenz.

Der frühe Kinetiktest ist entscheidend: Für das Produkt eines freien
masselosen selbstdualen Feldes und eines freien masselosen Weylfeldes
verschwindet die Ableitungsauslese exakt, auch bei nichtkollinearen
Impulsen. Ein sichtbarer Kanal benötigt eine wirklich berechnete andere
Kinetik oder Feldzuordnung.

Eine wörtliche skalare Zuordnung von b†ff mit linkem f macht b† links
und b rechts. Dann ist D=bf† gleichhändig rechts, χ†=b†ηf†
gemischthändig. Die Ladungen −1 und +3 bleiben verschieden. Native
Fockmoden sind noch keine vollständigen lokalen Felder.

\subsection{Unverändert offen und nächste
Entscheidung}\label{unveruxe4ndert-offen-und-nuxe4chste-entscheidung}

\textbf{Alle T1-T8-Gesamtpflichten bleiben offen.} Keine neue globale
Grundzustandswahl, 3+1D-Raumzeit, chirales Maß oder dynamische
Gravitation wurde bewiesen. Die v1.6.8-Polschranken wurden nicht
verbessert und dürfen nicht auf die neuen Kompositionsmodelle übertragen
werden.

Die nächsten Arbeiten sind jetzt scharf prüfbar: (1) Quellobjekte
liefern die Einbettungen und Überlappung ohne Fit; (2) ein gemeinsamer
autonomer Controller liefert den Impuls und Zeiger samt Energiebilanz;
(3) derselbe Feldvertrag liefert ein nichtnulles, vertexkonsistentes
Matrixelement. Erst anschließend trägt eine räumliche Skalierung.

\textbf{Prüfung:} 514 eigene exakt arbeitende Bedingungen und vier
numerische Kontrollen je normalem/optimiertem Lauf, identische
Ergebnisse. Separat 17 native und 1073 P0/P1-Quellguards. Vollständige
Beweise, Annahmen und historische Hauptfassung stehen im gleichzeitig
aktualisierten Hauptdokument und Prüfpaket.

\end{document}
